\documentclass[11pt]{article}
\usepackage[margin=21mm]{geometry}
\usepackage{amsmath,amssymb,amsthm,hyperref}
\hypersetup{hidelinks}
\newtheorem{lemma}{Lemma}
\newtheorem{proposition}[lemma]{Proposition}
\newcommand{\diag}{\operatorname{diag}}
\title{Half-line completion of the canonical Bell slack}
\author{}\date{}
\begin{document}\maketitle
\vspace{-2em}

A quadratic selector on each Bob half-line certifies a rational canonical
source that no global polynomial selector can certify. The additional
freedom has a simple Gram representation: a quadratic matrix positive on
$v\ge0$ is a linear matrix square plus $v$ times a constant positive matrix.
Applied to the canonical completion, this gives four constant positive
semidefinite blocks and a closed source cone. We give the complete affine
pencil and a rational certificate exhibiting the gain.

\section{The canonical completion}

Let $c=(d,\lambda,n,m,C,A,D,\alpha,\beta)\ge0$ and $D_v=1+v^2$.
The canonical state change collects the source into four quadratic matrices:
\begin{align*}
 P(v)&=\begin{pmatrix}4mv^2&-4mv\\-4mv&4m+(4\beta+2d)D_v\end{pmatrix},&
 N(v)&=\begin{pmatrix}4CD_v+4D+2d&(4D+2d)v\\(4D+2d)v&(4D+2d)v^2\end{pmatrix},\\
 E(v)&=\begin{pmatrix}4AD_v+(4\alpha+2d)v^2&-(4\alpha+2d)v\\
              -(4\alpha+2d)v&4\alpha-2dv^2\end{pmatrix},&
 K(v)&=4\lambda\begin{pmatrix}1&v\\v&v^2\end{pmatrix}
                  +4nD_v\begin{pmatrix}0&0\\0&1\end{pmatrix}.
\end{align*}
For $U=u^2$ and independent $a,b\in\mathbb R^2$, the reduced state form is
\begin{equation}
 \Phi_c=(1+U)a^TPa+U(1+U)b^TNb
             +U(a-b)^TE(a-b)+(a+Ub)^TK(a+Ub).             \label{eq:normal}
\end{equation}
The substitution
\[
 W=a_1+Ub_1,\quad X=a_2+Ub_2,\quad
 Y=u(a_1-b_1),\quad Z=u(a_2-b_2)
\]
gives the complete canonical parent $F_c=(1+u^2)\Phi_c$.
It is invertible when $u\ne0$, with determinant $u^2(1+u^2)^2$.
Thus positivity of this form covers all state directions; the finite
axis follows by polynomial continuity. Homogenizing the parent to
bidegree $(2,2)$ in the two projective angle pairs gives the infinite
endpoints. Real state positivity also covers complex states by separating
their real and imaginary parts.

A real $2\times2$ selector $J$ gives the two completion matrices
\begin{equation}
 G_J=\begin{pmatrix}P+K&J\\J^T&N+K\end{pmatrix},\qquad
 H_J=\begin{pmatrix}P+E&K-E-J\\K-E-J^T&N+E\end{pmatrix}.     \label{eq:GH}
\end{equation}
Their role is the exact cancellation
\begin{equation}
 \Phi_c=\binom a{Ub}^{\!T}G_J\binom a{Ub}
                  +U\binom ab^{\!T}H_J\binom ab.           \label{eq:split}
\end{equation}
We seek $J_+(v)=J_0+vJ_1+v^2J_2$ making both matrices positive
semidefinite for $v\ge0$. Their parity makes this sufficient for all Bob
angles: with $T=\diag(1,-1)$, each of $P,N,E,K$ satisfies
$M(-v)=TM(v)T$. For $v<0$ set $J_-(v)=TJ_+(-v)T$.
Congruence by $\diag(T,T)$ then restores both inequalities on that half-line.
At zero use $J_+(0)$; no continuity of the selector is needed.

\section{Four constant Gram blocks}

\begin{lemma}\label{lem:halfline}
A real symmetric $4\times4$ quadratic $M(v)=M_0+vM_1+v^2M_2$
is positive semidefinite on $v\ge0$ if and only if some unrestricted
$4\times4$ matrix $B$ satisfies
\[
 S=\begin{pmatrix}M_0&B\\B^T&M_2\end{pmatrix}\succeq0,
 \qquad T_M=M_1-B-B^T\succeq0.
\]
\end{lemma}
\begin{proof}
Sufficiency follows from $M(v)=[I\ vI]S[I\ vI]^T+vT_M$.
For necessity, $M(t^2)$ is positive on the real line and has degree at
most four. Matrix Fej\'er--Riesz gives polynomial factors of degree at most
two, also for singular matrices. Taking real and imaginary parts and
averaging under $t\mapsto-t$ gives
\[
 M(v)=(R_0+vR_2)^T(R_0+vR_2)+vR_1^TR_1.
\]
Take $S=[R_0\ R_2]^T[R_0\ R_2]$ and $T_M=R_1^TR_1$.
The degree bounds are sharp: $vI$ requires the half-line term and $v^2I$
requires a linear factor. No coefficient is divided out. The assertion
transfers to real closed fields because its finite Gram formulation is
first order. The real-line factorization used here is
\cite[Theorem 1.2]{Zalar}.
\end{proof}

Write $G_j,H_j$ for the coefficients of $v^j$ in \eqref{eq:GH}, with
$J=J_0+vJ_1+v^2J_2$. They are affine in $c$ and the selector entries.
Introducing two unrestricted $4\times4$ matrices $B_G,B_H$ gives the pencil
\begin{equation}
\boxed{\begin{aligned}
 S_G&=\begin{pmatrix}G_0&B_G\\B_G^T&G_2\end{pmatrix}\succeq0,&
 S_H&=\begin{pmatrix}H_0&B_H\\B_H^T&H_2\end{pmatrix}\succeq0,\\
 T_G&=G_1-B_G-B_G^T\succeq0,&
 T_H&=H_1-B_H-B_H^T\succeq0.
\end{aligned}}                                             \label{eq:pencil}
\end{equation}
Lemma~\ref{lem:halfline} proves both inclusions between this projection
and the quadratic half-line selector cone. It also retains every singular
case and every zero source face. There are $12+16+16=44$ auxiliary scalars
and no remaining coefficient equations. The block orders $8,8,4,4$ total
$24$; adding the nine scalar blocks $[c_i]$ gives total order $33$.

There is an explicit strictly positive dual certificate. Set
\[
 W_S=\begin{pmatrix}I_4&I_4/2\\I_4/2&I_4\end{pmatrix},\qquad W_T=I_4/2.
\]
Every auxiliary coefficient cancels in the exact identity
\begin{equation}
 \langle W_S,S_G+S_H\rangle+\langle W_T,T_G+T_H\rangle
      =16\sum_i c_i,                                      \label{eq:dual}
\end{equation}
where $\langle A,B\rangle=\operatorname{tr}(A^TB)$.
Indeed the $B$ traces cancel, the middle coefficients $G_1,H_1$ have zero
diagonal, and $\operatorname{tr}S_G+\operatorname{tr}S_H=16\sum c_i$.
Both weights have minimum eigenvalue $1/2$, so feasibility implies
\[
 \operatorname{tr}S_G+\operatorname{tr}S_H+
 \operatorname{tr}T_G+\operatorname{tr}T_H\le32\sum_i c_i.
\]
All $B_G,B_H,J_0,J_2$ entries are off-diagonal Gram entries.
For a $J_1$ entry at distinct $G_1$ indices $r,t$, combining the local
PSD entry bounds in $G_1=T_G+B_G+B_G^T$ gives
$|(J_1)_{ij}|\le(\operatorname{tr}T_G+\operatorname{tr}S_G)/2$.
Consequently every feasible auxiliary vector satisfies
\begin{equation}
 \|y\|_\infty\le16\sum_i c_i.                               \label{eq:bound}
\end{equation}
Bounded subsequences now prove the projection closed; replacing
$\sum c_i$ by $\sum|c_i|$ covers unrestricted source inputs.
The positive weights in \eqref{eq:dual} give strict feasibility of the
dual auxiliary equations. Fixed-source feasibility cannot be weakly
infeasible: an unbounded near-feasible sequence, normalized by its auxiliary
norm, would give feasible source-zero data with a nonzero auxiliary vector,
contradicting \eqref{eq:bound}; a bounded sequence gives a feasible limit.

\clearpage
\section{An exact source rescued by the half-line completion}

Take
\[
 c_* =\left(359,0,0,\frac{75}{16},60,416,
                   \frac{41}{2},\frac{831}{2},\frac{151}{2}\right).
\]
Dividing the source and every matrix below by $359$ gives $d=1$.
An exact selector on $v\ge0$ is
\begin{equation}
 J_+(v)=\frac1{64}\begin{pmatrix}
 3043v-2109v^2&3875v^2\\
 1403+23874v+2306v^2&815v+43455v^2
 \end{pmatrix}.                                           \label{eq:witnessJ}
\end{equation}
The two remaining auxiliary matrices are
\begin{align}
 B_G&=\frac1{256}\begin{pmatrix}
 0&0&0&0\\-4800&-48015&-27346&3260\\
 12172&186862&-21340&204800\\0&0&0&0
 \end{pmatrix},                                           \label{eq:witnessBG}\\
 B_H&=\frac1{18496}\begin{pmatrix}
 -6910952&240550&5277344&1292578\\
 -41117866&-4908665&35528504&2316913\\
 3043153&1489795&-7316324&276862\\
 36194938&-415582&-30578223&-3325523
 \end{pmatrix}.                                           \label{eq:witnessBH}
\end{align}
Together with \eqref{eq:GH} and \eqref{eq:pencil}, these rational matrices
specify the certificate completely. Their positive semidefiniteness is
verified by the exact elimination in the appendix. Thus \eqref{eq:split}
proves the full source valid for all states and all testing angles.

The difference from a global polynomial completion is essential.
If a polynomial $J(v)$ satisfies \eqref{eq:GH} on the real line, its entries
have degree at most two: each squared entry is bounded by a product of
two quadratic diagonal entries of $G_J$. Parity averaging makes its
diagonal entries even and its off-diagonal entries odd. On the face
$\lambda=n=0$, the zero diagonals of $G_J(0)$ and
$|J_{12}(v)|^2\le4m(4D+2d)v^4$ then reduce it to
\[
 J(v)=\begin{pmatrix}pv^2&0\\rv&sv^2\end{pmatrix}.
\]
For this source the exact contact is
$U=1/4$, $v=2$, $a=(17,2)^T$, $b=(12,-18)^T$.
The nonnegative terms in \eqref{eq:split} vanish there, forcing
$J(2)b=-4P(2)a$ and $J(2)^Ta=-N(2)b/4$.
These equations uniquely give $(p,r,s)=(-100,2000,1200)$.
At $v=1$ their $G_J$ principal submatrix on $(a_2,b_1,b_2)$ is
\[
 \begin{pmatrix}8235/4&2000&1200\\2000&1280&800\\1200&800&800\end{pmatrix},
 \qquad \det=-412640000<0.
\]
Consequently no global polynomial selector works for this source.
The reflected selector \eqref{eq:witnessJ} instead contains $|v|$ and
$v|v|$: its upper entries are
$(3043|v|-2109v^2)/64$ and $3875v|v|/64$; for $v\ne0$ the lower-left
entry also contains $1403\operatorname{sign}(v)/64$.

The result is an exact closed affine cone of quadratic half-line selectors,
strictly more effective than global polynomial selectors on this example.
Whether it contains every valid canonical source remains unresolved;
the order-$33$ pencil is therefore an inner certificate cone, not an
established affine representation of the full quantum set.

\clearpage
\appendix
\section{Exact positivity of the rational Gram matrices}

For the source above, $K=0$ and the three nonzero source matrices are
\begin{gather*}
 P=\begin{pmatrix}75v^2/4&-75v/4\\-75v/4&4155/4+1020v^2\end{pmatrix},\\
 N=\begin{pmatrix}1040+240v^2&800v\\800v&800v^2\end{pmatrix},\qquad
 E=\begin{pmatrix}1664+4044v^2&-2380v\\-2380v&1662-718v^2\end{pmatrix}.
\end{gather*}
Use \eqref{eq:pencil} with \eqref{eq:witnessJ}--\eqref{eq:witnessBH}.
Unpivoted symmetric elimination gives $M=LDL^T$ with unit diagonal in $L$.
At a positive residual pivot $d_i$, set $L_{ji}=R_{ji}/d_i$ and update
$R_{jk}\leftarrow R_{jk}-R_{ji}R_{ik}/d_i$ for $j,k>i$.
At every zero pivot the remaining row vanishes exactly, so it is left
unchanged. The following table lists all zero positions (numbered from one)
and the floors of the positive rational pivots, in order. Every positive
pivot lies in $[k,k+1)$ for its printed integer $k$.
\[
\begin{array}{c|c|l|c}
 M&\text{zero positions}&\text{floors of positive pivots}&\operatorname{rank}M\\\hline
 S_G&1,4,8&1038,\ 1039,\ 16,\ 382,\ 162&5\\
 S_H&8&1664,\ 2700,\ 1039,\ 639,\ 1562,\ 214,\ 67&7\\
 T_G&1,3,4&375&1\\
 T_H&4&747,\ 483,\ 430&3
\end{array}
\]
For example,
$T_G=(5335/128)(0,3,-2,0)^T(0,3,-2,0)$.
The printed rational inputs and this recurrence give a pen-and-paper
positivity test including all singular directions. The accompanying exact
witness file also supplies every entry of $L,D$; the rational checker
recomputes them from these inputs and verifies $M=LDL^T$.

\begin{thebibliography}{1}
\bibitem{Zalar} A. Zalar, \emph{Matrix Fej\'er--Riesz theorem with gaps},
Theorem 1.2, \href{https://arxiv.org/abs/1503.06034}{arXiv:1503.06034}.
\end{thebibliography}
\end{document}
